\documentclass[../main.tex]{subfiles}
\begin{document}
%==========================================TIÊU ĐỀ TẬP========================================
\begin{table}[h]
    \begin{adjustwidth}{-1em}{}
        \begin{tabular}{>{\centering\arraybackslash}p{6cm}|>{\raggedright\arraybackslash}p{12.5cm}}
            \multirow{5}{6cm}
            % Cột bên trái: ÉPISODE và số tập
            {\rainbowcolor
                {\\[-40pt]
                    \flaregothic\fontsize{20pt}{20pt}\selectfont \centering
                    \color{eptype}JOUR\color{black}\\[12pt]
                    \hbrk\fontsize{72pt}{72pt}\selectfont
                    \color{epnum}6
                    \\[18pt]
                    % \flaregothic\fontsize{14pt}{14pt}\selectfont
                    % \color{epattr}BIENVENUE, \uline{\color{Banpire} Mel} !
                    % \flaregothic\fontsize{18pt}{18pt}\selectfont
                    % \color{epattr}ÉPISODE PERDU\\
                    \unrainbowcolor}
                \color{black}}
             &
            % Dòng thứ nhất: tên tập tiếng Nhật
            {\fotskip\fontsize{18pt}{18pt}\selectfont きょうのたまき
            }                                                                                     \\[6pt]
            % Dòng thứ hai: tên tập tiếng Anh
             & {\cafeteria\fontsize{18pt}{18pt}\selectfont Fumino Tamaki.}                                                                       \\[4pt]
            % Dòng thứ ba: tên tập tiếng Việt
             & {\vnmsans\fontsize{18pt}{18pt}\selectfont Câu chuyện thường nhật về nàng mèo hoang nghịch ngợm và những nữ sinh trung học}        \\[8pt]
            % Dòng thứ tư: Nhân vật xuất hiện
             & {\sen{\fontsize{14pt}{14pt}\selectfont Track used: Out of Control Era (from Professional Block) (Supercharge - MDK \& meganeko)}} \\[6pt]
            % Dòng cuối cùng: Thời gian diễn ra của tập (thường sẽ là ngày phát hành)
             & {\continuum\fontsize{16pt}{16pt}\selectfont Ngày phát hành: 24/07/2021}                                                           \\[18pt]
        \end{tabular}
    \end{adjustwidth}
\end{table}
%=========================================TRUYỆN CHÍNH========================================
\color{black}

\tc{\uline{Quy ước trong phần bình luận:}}
\tl{Đây là lời thoại của Fumi.}
\tr{Đây là lời của Kuon, sẽ căn lề phải.}
\tr{[Và đây là lời của Suzuki, sẽ căn lề phải, đóng khung trong ngoặc vuông.]}

\il

\pov{Bình luận của Fumi, Hikawa Kuon và Hikawa Suzuki.}

\tc{Chào mừng quý vị đến với chương trình ``Nhật ký thường nhật của Tamaki.''}

Một ngày dịu dàng đến bất ngờ tại văn phòng, khi ánh nắng rọi vào ô cửa kính, chúng tôi bắt gặp một khung cảnh quen thuộc. Trên chiếc ghế sô-pha mềm mại, là một cô mèo đang lăn lộn ở kia. Đó là Fumino Tamaki.

\img{nj-6a}

``Meo\~{}''

\tr{[Như bao lần khác, cô mèo hoang nhỏ bé của nhà Cầu Vồng lại đang say giấc. Cứ như thể thời gian đối với em ấy chẳng có chút ràng buộc nào vậy.]}

\tr{Anh phải công nhận là Tamaki-san rất đáng yêu. Tóc ngắn, hơi nâu, phong cách giản dị\dots\ Làm anh nhớ đến Okayu-chan bên phía \hll. Cũng là một cô mèo mê ngủ đấy.}

Nhưng Tamaki không chỉ biết ngủ đâu. Em ấy rất hiếu động. Dù sao thì\dots\ em ấy vẫn là một chú mèo hoang.

\begin{wrapfigure}[27]{l}{8cm}
    \includegraphics[width=8cm]{../../../../Image/Book?/tamaki_model.jpg}
\end{wrapfigure}

\tr{[Ý chị là `noraneko' đúng không ạ? Thuật ngữ đó không chỉ dùng để chỉ loài mèo thật, mà còn để nói về cá tính hoang dã, tự do\dots\ Có vẻ rất hợp với Tamaki-nee đấy ạ.]}

Chính xác. Ở công ty, ai cũng gọi em ấy là ``mèo hoang'' không phải vô cớ đâu.

\chrint

Có lời kể rằng, Fumino Tamaki (\ruby{文}{ふみ}\ruby{野}{の}\ruby{環}{たまき}, \ruby{Phư-mi}{Văn}-\ruby{nô}{Dã}\ \ruby{Ta-ma-ki}{Hoàn}) từng là một cô mèo hoang đích thực. Vì thế mà chẳng ai biết được em ấy sinh ngày nào, thậm chí tuổi thật là bao nhiêu.

\tr{Vậy nên anh mới thấy có cái ``truyền thuyết đô thị'' rằng sinh nhật của em ấy là\dots\ ``mỗi ngày''. Cũng hợp lý nhỉ?}

*cười khẽ* Có lẽ đúng là như vậy. Nhưng em ấy không chỉ nằm không rỗi hơi đâu. Thỉnh thoảng Tamaki còn thức dậy đi lại xung quanh để canh chừng khu vực đấy.

\tr{[Vậy chị ấy là bảo vệ ngầm cho công ty ạ?]}

Không, chỉ là bảo vệ cho lãnh địa của mình thôi. Bất cứ ai dám mon men lại gần cũng đều sẽ bị em ấy tấn công không thương tiếc.

\tr{\dots Nghe giống như một con mèo hàng thật nhỉ.}

Nhưng nếu một ai đó không có ý định chiếm lãnh địa, em ấy sẽ trở nên ngoan hiền và thậm chí còn xin một lon cá nữa.

\tr{Ồ.}

\tr{[Chị ấy cũng khá nổi tiếng trên các nền tảng phát trực tuyến đúng không ạ?]}

Đúng vậy. Tamaki là một cô mèo rất thú vị. Mỗi buổi phát sóng của em ấy đều khá là\dots\ nói sao nhỉ, độc đáo.

\tr{Anh còn nghe nói rằng em ấy từng làm một buổi stream game kinh dị, nhưng thay vì la hét thì lại ngủ quên giữa chừng. Khán giả không biết nên cười hay nên lo.}

*cười* Đó là chuyện thường ngày mà. Có những lúc em ấy phát sóng rất bình thường, nhưng đột nhiên lại có những khoảnh khắc kỳ quặc đến mức người xem không biết đó là thật hay đùa nữa.

\tr{[Em còn nhớ vụ động đất năm 2021 ấy ạ. Em có đọc tin trên X, thấy chị ấy nói rằng đang ngủ thì bị lò vi sóng rơi trúng tay, phải đi cấp cứu luôn.]}

Ừm, may là không có gì nghiêm trọng. Nhưng đó cũng là một ví dụ điển hình về cuộc sống thường ngày của Tamaki đấy. Mọi thứ cứ\dots\ bất ngờ kiểu gì ấy.

\tr{[Như một chú mèo thật vậy - lúc nào cũng tạo ra những tình huống không ai ngờ tới.]}

Nhưng chính vì thế mà mọi người mới yêu quý em ấy. Tamaki-san có một sức hút rất riêng - vừa đáng yêu, vừa hài hước, lại còn có chút gì đó\dots\ không thể đoán trước được. Dù có kỳ quặc thế nào đi nữa thì Tamaki vẫn là Tamaki của chúng ta. Một cô mèo có mái tóc màu nâu sẫm dài đến bả vai, đôi mắt màu lục bảo, và đôi tai mèo dựng đứng cùng cái đuôi sọc nâu-trắng phía sau.

\tr{Tò mò một chút, bộ trang phục kia mà em ấy đang mặc là sao vậy?}

À, bộ áo hoodie mang sọc nâu, chiếc áo đen mặc ở bên trong, váy ngắn màu nâu, đi quần tất đen quá đầu gối và giày thể thao đỏ đúng không? Từ lúc em ấy được mang về thì đã thế rồi. Thậm chí tấm bảng ``làm ơn cho tôi thức ăn đóng hộp cho mèo'' đeo trước ngực cũng đã ở đó từ bao giờ.

\tr{[Vậy mấy thứ cài tóc trên đầu chị ấy thì sao ạ?]}

Cái kẹp tóc hình cá màu xanh ở mái kia là được chị tặng. Em ấy thực sự mê tít nó lắm.

\tr{[Vậy, ai là người đầu tiên ``nhặt'' chị ấy về ạ?]}

Không ai dám chắc, nhưng lời đồn kể rằng\dots\ có lẽ là một nhóm người đã vô tình gặp em ấy, rồi mang về công ty.

\tr{Anh đoán là\dots\ ba cô nàng kia phải không? Ba cô gái mặc bộ đồ trung học tám chuyện với nhau đấy.}

Bingo. Không sai một chút nào.

\tr{Anh dám cá là khi đó Tamaki-san vùng vẫy ghê gớm lắm. Kiểu, em ấy vừa quẫy vừa kêu ``Thả tôi ra! Thả tôi ra, nyaaa!'' như thể đang bị bắt cóc vậy.}

\begin{figure}[H]
    \centering
    \begin{subfigure}[t]{0.45\textwidth}
        \centering
        \includegraphics[width=8cm]{../../../../Image/Book?/nj-6b1.png}
    \end{subfigure}
    \quad
    \begin{subfigure}[t]{0.45\textwidth}
        \centering
        \includegraphics[width=8cm]{../../../../Image/Book?/nj-6b2.png}
    \end{subfigure}
    \quad
    \begin{subfigure}[t]{0.45\textwidth}
        \centering
        \includegraphics[width=8cm]{../../../../Image/Book?/nj-6b3.png}
    \end{subfigure}
    \caption{Quá trình mà nhóm JK-gumi (Câu lạc bộ Nữ sinh trung học) bắt Tamaki về.}
\end{figure}

*cười nhẹ* Đúng như anh nghĩ đấy.

\tr{[Tội nghiệp chị ấy ghê\dots]}

\chrint

Nhưng giờ đây, Tamaki đã có một mái ấm mới, ngay tại văn phòng Nijisanji này.

\tr{[Dù đã có chốn về, em thấy chị ấy vẫn giữ nguyên nét nghịch ngợm của một chú mèo hoang. Nhưng chính điều đó khiến Tamaki-nee trở nên đặc biệt, đúng không ạ?]}

Chính xác.

\tr{Này nhé, anh mới thấy em ấy vừa nhảy lên bàn rồi đấy. Không ai cản nổi đâu.}

Bản năng thôi. Với Tamaki, không ai có thể đoán trước được hành động tiếp theo.

Mỗi người hiện giờ đang thưởng thức một chiếc dorayaki kèm một ly trà xanh ấm. Nhìn thế này đúng là\dots\ một buổi chiều ấm cúng.

\tr{[Khung cảnh yên bình quá nhỉ. Em phải thừa nhận là\dots\ kể cả khi Tamaki-nee leo lên bàn hay làm điều gì bất ngờ, thì sự hiện diện của chị ấy vẫn khiến người khác thấy bất ngờ.]}

Mà nhắc mới nhớ\dots\ chúng ta cũng nên giới thiệu sơ qua về ba cô gái mặc đồng phục nữ sinh đang ngồi kia, đúng không?

\chrint

\tc{\uline{(Cảnh quay tạm dừng, máy quay lia chậm qua từng thành viên nhóm JK-gumi)}}

\begin{wrapfigure}[27]{r}{8cm}
    \includegraphics[width=8cm]{../../../../Image/Book?/kaede_model.png}
\end{wrapfigure}

\tr{Ý hay đấy. *e hèm* Trước tiên, cô nàng vừa mới nhìn Tamaki-chan với sự buồn phiền. Higuchi Kaede (\ruby{樋}{ひ}\ruby{口}{ぐち}\ruby{楓}{かなえ}, \ruby{Hi}{Thông}-\ruby{gư-chi}{Khẩu}\ \ruby{Ca-ê-đê}{Phong}), trưởng nhóm ``Câu lạc bộ Nữ sinh Trung học.'' Cô ấy là một nữ sinh trung học năm hai ở vùng Kansai ảo, gia nhập vào gia đình Cầu Vồng vì muốn mọi người biết nhiều hơn về những thứ mình thích.}

Kaede-chan có điểm nổi bật nhất nằm ở mái tóc cột đuôi ngựa dài tới quá lưng, đôi mắt tím mê hoặc và nốt ruồi dưới mắt phải, và đó cũng là điểm khiến mọi người thích nàng.

\tr{[Vậy chị ấy có điểm gì thu hút nữa không ạ?]}

Cô ấy nổi tiếng với điệu cười ``cao quý'' và cách mà cô ấy nói ``Nande!! (Tại sao chứ!)'' hài hước của mình.

\tr{Anh dám chắc là ở \hll\ cũng có ai đó có cách nói ``Nande\~{}'' như thế. Có lẽ vậy.}

Và nếu hai người không biết, Kaede-chan có cách nói giống một nữ sinh nhất trong nhóm, mang sắc thái trẻ trung nhất, tới độ mà đôi khi cô được gọi là ``Nữ sinh trung học Nijisanji'' đấy.

\tr{[Và cũng là người cao nhất trong cả ba đấy ạ.]}

Nghe nói cô ấy cũng ở trong ban nhạc của trường nữa.

\tr{[Vậy chị ấy chơi gì ạ? Trống? Sáo? Hay là-]}

Kèn trumpet.

\tr{\dots Hả?}

Cô ấy vốn dĩ đã thích chơi nhạc cụ đó từ lúc còn nhỏ, và hiện đang là thành viên của một ban nhạc kèn đồng và diễu hành nữa.

\tr{Dù là ở nhà hay lên công ty, cô ấy luôn diện bộ đồng phục của mình nhỉ. Áo sơ-mi trắng, cà-vạt tím, áo vest nâu, áo khoác ngoài màu đen đi cùng với chiếc váy họa tiết ca-rô màu đen xếp ly, đúng là bộ đồng phục tiêu chuẩn của một nữ sinh trung học.}

\tr{[Đi cùng với đó là tất cao đến đầu gối màu trắng và giày lười đen nữa. Cũng là một bộ hoàn hảo cho đồng phục nữ sinh ạ.]}

\begin{wrapfigure}[27]{l}{8cm}
    \includegraphics[width=8cm]{../../../../Image/Book?/mito_model.png}
\end{wrapfigure}

Tiếp đến là Tsukino Mito (\ruby{月}{つき}\ruby{ノ}{の}\ruby{美}{み}\ruby{兎}{と}, \ruby{Xư-ki}{Nguyệt}-nô \ruby{Mi}{Mỹ}-\ruby{tô}{Thỏ}), một nữ sinh năm hai giữ chức lớp trưởng gương mẫu, nhưng có đôi lúc tsundere và đôi lúc cũng trái ngược hoàn toàn so với những gì mình định hướng. Cô ấy được coi là ``lá cờ tiên phong cho Nijisanji chúng em đấy.''

\tr{Đối trọng với Tokino Sora-chan ở bên bọn anh nhỉ?}

Có thể nói là như vậy.

\tr{[Mà Fumi-nee, chị nói Mito-nee là ``trái ngược hoàn toàn so với những gì mình định hướng'' là sao ạ?]}

Thì, những lần cô ấy chơi ``kusoge'' lố bịch hoặc mấy trò dạng flash game đơn giản, hay những lựa chọn phim kỳ quặc, và những điều ngớ ngẩn cô ấy từng làm đó.

\tr{[\dots Thật ạ?]}

Ừ. Còn cái tính tsundere mà chị vừa nói, thì đó là lúc khi cô cố gắng trở thành một lớp trưởng nghiêm túc. Mặc dù đã cố gắng hết sức, Mito vẫn cảm thấy những nỗ lực của mình là vô ích. Cô ấy thường tự hỏi liệu mình có nói quá nhiều sau khi phát biểu không và cảm thấy chán nản sau đó.

\tr{Cái đó gọi là không tự tin rồi còn gì. Mà điều ngớ ngẩn em nói tới là sao?}

Có lần trong khi chơi A.Crossing, cô ấy đã biến buổi stream của mình thành\dots\ một bộ phim truyền hình dài tập với những mối quan hệ, sự phản bội và nỗi đau khổ của các thành viên với nhau. À, còn nữa, Mito cũng muốn mọi người vẽ nhiều fanart và manga dâm đãng về bản thân nữa.

\tr{\dots Giống như một người có sở thích từ ai đó đội mũ thuyền trưởng lai với một cô chó nào đó ở bên anh nhỉ.}

Cô ấy mà nghe được anh nói là kiểu gì chả hét lên: ``\textbf{K-Không phải như anh nghĩ đâu!!}'' như thế đấy.

Dù gì đi nữa, cô ấy vẫn trở thành một trong những người mà mọi người ở Nijisanji quý trọng và yêu mến. Và vì cô nàng cũng là một người trong JK-gumi\dots

\begin{wrapfigure}[29]{r}{6.5cm}
    \includegraphics[width=6.5cm]{../../../../Image/Book?/rin_model.png}
\end{wrapfigure}

\tr{[\dots nên em đoán chị ấy cũng diện bộ đồng phục trường học rồi. Áo khoác đen ngoài, áo len màu nâu bên trong, với áo sơ-mi trắng, đeo nơ hồng trước ngực và váy xếp ly - em nghĩ là cùng loại với Kaede-nee, cùng với tất trắng dài tới tùi mang hoa văn thỏ và giày lười đen. Em nghĩ hai chị ấy học cùng trường\dots]}

\tr{Ngoài ra, Mito còn có mái tóc đen dài thẳng với hai chiếc kẹp màu bạc và hồng đậm, em ấy có buộc tóc quấn quanh đầu và có đôi mắt tím nữa.}

Và cuối cùng là ai nào, Suzuki-chan?

\tr{[Shizuka Rin (\ruby{静}{しずか}\ruby{凛}{りん}, \ruby{Si-du-ca}{Tịnh}\ \ruby{Rin}{Lẫm}), người lớn tuổi nhất trong bộ ba chị em nữ sinh này. Chị ấy học ở năm ba và là người có tính cách trưởng thành nhất trong ba người họ, theo như những gì em đọc ở đây.]}

Ồ, cô nàng có tính cách như một người chị cả ha\~{} Tính cách nhẹ nhàng, luôn được tin tưởng và yêu mến, chả phải quá đúng sao?

\tr{Dám chắc là em ấy được nhiều người thích lắm.}

Đương nhiên. Kể cả là bọn con trai. Nhưng mà cô ấy toàn tìm cách lảng tránh họ đi.

\tr{\dots Sao nghe hao hao giống như anh vậy, chả hiểu sao mà anh có cảm giác mình là cục nam châm thu hút các thành viên bên \hll được\dots}

Chắc là vì anh có hội tụ đủ yếu tố như Rin sao?

\tr{[Onii-sama cũng là một người tốt bụng và giúp đỡ mọi người ở công ty mà. Anh ấy cũng là người giúp em về được với Onee-sama, và cũng là người thuyết phục bọn em đi theo anh ấy đấy!]}

\tr{\dots Anh không nghĩ là mình giỏi giang đến thế đâu. *e hèm* Ngoài những tính cách ở trên, Mito-san còn có những đặc điểm nào nữa?}

Cô ấy nổi tiếng là một người stream rất dài. Một buổi stream của cô có thể kéo dài tới hơn mười tiếng, kéo dài từ tối hôm trước sang sáng hôm sau mà.

\tr{\dots Khiếp vãi. Y như Koro-chan nhà anh luôn.}

*cười nhẹ* Cô ấy là một game thủ chính hiệu mà. Thậm chí cô ấy còn cố gắng lấy hết tất cả các thành tựu của một game hay cố gắng tới một cái kết trong game đó, dù rằng nó bất khả thi nếu chỉ ngồi lần chơi.

\tr{Nói như thế tức là cô ấy chơi game giỏi lắm nhỉ?}

Cũng không hẳn. Tùy từng game nữa. Chỉ là một người chơi game bình thường thôi.

\tr{[Vậy Rin-nee cũng học cùng trường vơi Mito-nee và Kaede-nee chứ ạ?]}

Không. Bộ đồng phục của cô ấy khác biệt hẳn mà. Bộ đồ cô ấy gồm áo khoác ngoài màu xanh đậm, váy xếp ly không có hoạt tiết, áo sơ-mi trắng có cổ, áo len xám và đeo nơ tím hồng trước ngực. Cô ấy mặc quần tất và giày lười đen.

\tr{Hèn gì mà trông cô ấy khác hẳn so với hai người kia luôn.}

\chrint

\img{nj-6c}

\tr{[Vậy thì\dots\ ta nên tiếp tục thôi nhỉ, Onii-sama?]}

Được chứ. Giờ đến đoạn nào rồi nhỉ\dots\ À-rá?

\tr{Có vẻ như nàng Mèo hoang của chúng ta đang chui xuống gầm bàn đấy.}

\tr{[\dots Mặt của chị ấy còn đỏ ửng lên kìa. Ta vẫn không thể biết được Tamaki-nee đang nghĩ gì.]}

Như tôi đã nói từ đầu, Tamaki là một cô mèo rất khó hiểu. Em ấy nghịch ngợm, có phần vô tư quá mức. Nhưng đó cũng là một phần sức hút đặc biệt của em ấy.

\tr{Vậy nếu giờ em ấy bám lấy chân Rin-chan rồi ``gặm gặm'' như đang xử lý một cái xương, thì cũng là sức hút luôn à?}

\tr{[*cười* Ý anh là những gì mà chị ấy đang làm lúc này ạ?]}

\tr{Ừ. Cứ nhìn mà xem.}

\img{nj-6d}

Nếu là Rin thì chắc chẳng vấn đề gì đâu. Nhưng nếu là Kaede hay Mito thì\dots

\tr{Nhìn cái cách hai cô nàng kia đuổi theo Tamaki là hiểu ngay mức độ ``báo động'' rồi.}

Có người từng nói, chính sự phiền toái ấy lại khiến người ta không thể ghét Tamaki được.

\tr{[Giống như việc nuôi một chú mèo thật vậy. Có thể người ta sẽ phát điên vì những trò nghịch ngợm của chúng, hay việc làm bừa bãi khắp nhà\dots\ nhưng rồi lại chẳng nỡ mắng.]}

\tr{Chuẩn không cần chỉnh.}

\img{nj-6e1}

\dots Khoan, hình như Tamaki lại đang ngó nghiêng gì đó. Nhất là với đĩa bánh rán của Kanae kìa.

\tr{[Đúng như là thế thật. Chiếc dorayaki bên kia\dots\ có vẻ đang lọt vào tầm ngắm rồi kìa.]}

\tr{Anh sẽ không bất ngờ nếu em ấy lao thẳng tới và chôm mất.}

\begin{figure}[H]
    \centering
    \begin{subfigure}[t]{0.45\textwidth}
        \centering
        \includegraphics[width=8cm]{../../../../Image/Book?/nj-6e2.png}
        \caption{Tamaki nhảy lên, mắt sắc lẹm hướng tới cái bánh\dots}
    \end{subfigure}
    \quad
    \begin{subfigure}[t]{0.45\textwidth}
        \centering
        \animategraphics[autoplay,loop,width=8cm]{30}{../../../../Image/Book?/nj-6e3/nj-6e3.}{1}{47}
        \caption{\dots ngoạm lấy chiếc bánh nhanh như cắt, không để bất cứ ai trong ba cô gái phản ứng kịp\dots}
    \end{subfigure}
    \quad
    \begin{subfigure}[t]{0.45\textwidth}
        \centering
        \includegraphics[width=8cm]{../../../../Image/Book?/nj-6e4.png}
        \caption{\dots cuối cùng làm Kaede phát khóc (đúng nghĩa đen). Ta có thể nghe thấy cô ấy thốt lên: ``\textbf{Tại sao chứ!??}''}
    \end{subfigure}
\end{figure}

Trời ơi\dots\ Em ấy vừa khiến một cô gái bật khóc rồi kìa!

\tr{Anh đoán được mà. Động tác nhanh đến mức không ai phản ứng kịp.}

\tr{[Kaede-nee chỉ còn biết trân trân nhìn theo\dots\ như thể không tin nổi vào mắt mình.]}

\tr{*thở dài nhẹ* Người mất bánh thì khóc ròng, còn kẻ cướp bánh thì đang thưởng thức trong an nhà\dots\ đúng là nghịch lý của đời.}

Không biết rồi Tamaki có ổn không nữa\dots

\tr{Em đang lo là em ấy sẽ bị ``trả đũa'' à?}

\tr{[Fumi-nee cũng có ý của chị ấy đấy. Một cách nào đó thì\dots\ Tamaki-nee sắp gặp rắc rối rồi.]}

\tr{Ý em là kiểu ``Gòi, gòi! Mày xong rồi!'' phải không?}

\dots Cái cách anh nói nghe hơi khó hiểu, nhưng đúng tinh thần đấy. Sắp có chuyện lớn rồi.

\tc{\uline{Từ đằng sau Tamaki đang ăn chiếc bánh ngon lành, một thứ lửa bừng lên từ phía sau. Kaede tức tới mức mất hết con ngươi, tay lăm lăm cây chày gỗ đóng rất nhiều đinh}}

``\dots Grừ\dots''

``Meo?''

\img{nj-6f}

\dots Ôi trời. Kaede-chan thực sự giận rồi. Nhìn khí thế kia mà xem\dots

\tr{[Chị ấy lấy đâu ra cây gậy bóng chày đó thế!? Trông Kanae-nee không khác kẻ man rợ!]}

*RẦM! RẦM! RẦM!*

*VÚT!* *VÚT!*

\tc{\uline{Kanae xả hết cơn giận, cố gắng đập lấy cô mèo ăn vụng kia, cùng với những mảnh bàn ghế bị bay tứ tung.}}

Tôi có thể nghe thấy tiếng hét của Kaede\dots\ như thể bao nhiêu kiềm nén trong cô ấy vừa bùng nổ.

\tr{Em ấy quyết ăn thua đủ với Tamaki-chan luôn rồi\dots\ Văn phòng bị đập tan nát cả rồi còn đâu!}

\tr{[Hai cô bạn của chị ấy thậm chí còn\dots\ nắm chặt tay nhau vì sợ kia kìa.]}

Tama-chan\dots\ Em ấy vẫn ổn chứ? À, khoan\dots

\tc{\uline{Bụi khói tan dần, lộ ra cảnh tan hoang của văn phòng như vừa xảy ra một trận cướp phá, với Kanae đứng giữa phòng, thở hồng hộc, còn Mito và Rin ở đằng sau, run rẩy ôm chặt tay nhau.}}

\img{nj-6g}

\tr{[Giờ thì văn phòng của Nijisanji không khác gì sau bão. Y như ở căn hộ cũ nơi mà em với Onee-sama từng ở.]}

\tr{Làm anh nhớ lại những ngày đầu ở \hll\dots Ayame-chan hay Robo-PON cũng từng phá tan văn phòng y như thế này.}

\tr{[Nhưng ít nhất thì Tamaki-nee vẫn ổn\dots\ Đúng không?]}

Chị tin là vậy. À, xem kìa.

\tc{\uline{Mặc cho căn phòng giờ bừa bộn và đầy rẫy đống đổ nát, Tamaki - con mèo hoang của nhà Cầu Vồng - vẫn cuộn tròn lại và nằm một giấc ngủ ngon lành.}}

\img{nj-6h}

\tr{[Chị ấy lại\dots\ nằm xuống như chẳng có gì xảy ra.]}

\tr{Không thể tin nổi. Mọi chuyện vừa rồi như thể không liên quan gì đến em ấy vậy.}

Đó chính là Tamaki. Luôn sống trong thế giới riêng, luôn vô tư lạ thường\dots

\tc{Quả thật, Fumino Tamaki là một cô nàng mèo hoang nghịch ngợm.}
%==========================================TRUYỆN PHỤ=========================================
\omk

\tc{\uline{(Kết thúc phần bình luận của ba người.\\Văn bản sẽ chuyển về dạng \LaTeX\ bình thường.)}}

\pov{-}

\begin{wrapfigure}[29]{l}{11cm}
    \includegraphics[width=11cm]{../../../../Image/Book?/fumi_model.jpg}
\end{wrapfigure}

Fumi, Kuon và Suzuki vừa hoàn tất phần thuyết minh cho tập mới nhất của chương trình ``Tamaki thường nhật'' một chuỗi sê-ri mang tính tài liệu, do chính nàng Cáo thần khởi xướng và sản xuất.

\chrint

Fumi (フミ, Phư-mi) là thành viên gia nhập vào Nijisanji năm 2019. Cô là một vị thần Cáo được thờ phụng trong một ngôi đền Nhật Bản, người đã sống ẩn dật trong khoảng năm mươi năm. Chính bởi sống trong đó một thời gian rất dài, đôi khi cô có thể sử dụng những tiếng lóng hoặc cách suy nghĩ có phần lạc hậu.

Fumi thường tạo ấn tượng là một người chín chắn và trưởng thành thông qua cách nói chuyện và vẻ ngoài của mình. Tuy nhiên, đằng sau vẻ ngoài ấy lại là một tâm hồn tinh nghịch và đôi khi khá lạ lùng. Những buổi stream của cô thường có những khoảnh khắc khó hiểu, như việc ngồi đếm từng sợi tóc trên đầu hay tỉ mẩn tách từng thành phần trong gói gia vị rắc cơm.

Dù hồ sơ chính thức có thể không đề cập, nhưng Fumi lại là một người cuồng nhiệt của thể loại anime khoa học viễn tưởng. Cô đặc biệt say mê sê-ri Gundam và thường xuyên lên sóng để lắp ráp các mô hình gunpla trước sự theo dõi của người xem.

Là một thành viên tích cực của gia đình Nijisanji, Fumi luôn quan tâm theo dõi mọi hoạt động của các đồng nghiệp, kể cả những chi nhánh ở nước ngoài. Đặc biệt, cô có mối quan hệ gần gũi với Fumino Tamaki - người mà cô thường xuyên trò chuyện và bình luận về những hành động kỳ quặc của cô nàng mèo hoang này.

Fumi là một thần Cáo, nên dĩ nhiên cô có đôi tai cáo màu nâu và một cái đuôi rất lớn ở đằng sau, kèm với đôi mắt vàng, chỏm tóc ngố và một vài phụ kiện đính trên tóc và trên tai. Cô có mái tóc nâu dài tới hông, được tết thành bím. Bộ trang phục của cô là một bộ kimono trắng truyền thống thường được thấy ở các vu nữ trong chùa: váy hakama đỏ, tay áo tách rời, buộc bằng dây obi đỏ, tất tabi trắng, dép quai đỏ và đeo vòng cổ đỏ. Quanh thân cô còn có thêm hai sợi tơ rua khắc các hình màu vàng trên đó.

\chrint

``Vậy\dots'' Kuon liếc nhìn quanh khung cảnh tan hoang của văn phòng từ cửa sổ văn phòng bên cạnh, gãi đầu thở dài, ``đống đổ nát này chúng ta sẽ xử lý thế nào đây?''

Suzuki chắp tay sau lưng, cười trừ: ``E là\dots\ sẽ khá tốn sức đấy. Có thể mệt tới mức phải xin nghỉ vài hôm mất.''

``Ôi trời, có gì đâu,'' Fumi khẽ cười, giọng cô pha chút lả lơi. ``Cứ để ba cô gái nữ sinh kia tự lo đi. Dù sao thì Kanae cũng nên có trách nhiệm đứng ra giải quyết chuyện này. Giống như\dots\ đếm tóc trên đầu vậy.''

``Chị thực sự thấy hai việc đó giống nhau ở điểm nào vậy, Fumi-nee?'' cô bé khẽ nghiêng đầu, chân thành hỏi.

``Anh còn ngạc nhiên là những lời như thế lại có thể thốt ra từ một Thần Cáo đã tự nhốt mình trong đền suốt nửa thế kỷ,'' cậu Tổng chen ngang, tròn mắt nhìn.

Fumi chỉ khẽ nhún vai, mỉm cười như thể chưa từng nghe thấy lời bình đó. ``Có gì đâu mà ngạc nhiên thế\dots''

``Bọn anh/em không có khen em/chị đâu,'' cả hai anh em đồng thanh.

Một thoáng im lặng trôi qua, trước khi Kuon thở dài lần nữa: ``Thế ai sẽ là người xử lý chuyện làm dịu hai cô bé đang hoảng loạn ngoài kia? Mito-chan và Rin-chan nhìn như sắp khóc đến nơi rồi.''

\img{nj-6i}

Nàng Cáo thần thản nhiên trả lời: ``Họ sẽ lo liệu thôi.'' Tới mức này, cả hai anh em nhà Hikawa không còn lời nào để nói với cô nàng thảnh thơi tới độ chối bỏ trách nhiệm kia.

Ngay lúc đó, cánh cửa phòng thu mở ra, và bước vào không ai khác chính là ``thủ phạm'' khiến cả văn phòng rối tung rối mù ban nãy - Fumino Tamaki. Trông cô nàng như vừa mới ngủ trưa dậy, tóc tai rối bù, hai tai mèo lắc lư theo từng bước chân.

``Meo\~{} Em tới lồng lại giọng cho đoạn của em trong tập hôm nay đây\~{}'' cô nàng tuyên bố một cách tự hào, giơ một ngón tay lên trời như thể mình vừa làm điều gì đó vĩ đại.

Cô bạn Cáo của mình vẫn đang kiểm tra lại âm thanh hậu kỳ, chỉ nhẹ nhàng đáp: ``À, bọn chị thu xong hết rồi.''

``\textbf{Ể!!??}'' tiếng hét của nàng Mèo hoang vang lên như khi bị giẫm đuôi. ``Tại sao chứ!? Không có em thì làm sao có chất mèo hoang đích thực được?! Ai\dots\ ai lồng giọng cho em!?''

Fumi vẫn điềm tĩnh như cũ, quay màn hình lại cho cô nàng xem. Trên đó là một đoạn ngắn hậu trường: một con mèo thật - tai nhọn, lông mướt, đôi mắt lười biếng nhưng sắc sảo - đang nằm dài trong phòng thu âm, miệng ``meo\~{}'' đúng lúc và đúng tông đến kỳ lạ.

``\dots Đùa nhau à,'' Tamaki đứng đờ ra. ``Cái đó\dots\ là con mèo thật á!?''

``Chị bắt gặp nó ngoài bãi xe phía sau toà nhà,'' nàng Cáo mỉm cười. ``Nó có tố chất hơn em nhiều. Lồng giọng rất đúng chỗ, lại không phá phách nữa\~{}''

Nàng Mèo hoang quay phắt sang, và rồi, ánh mắt cô nàng chạm phải Kuon, người vừa thản nhiên nhấm nháp trà xanh, và Suzuki, người đang lật lại biên bản ghi hình với nụ cười mỉm.

``\dots Cả hai người cũng tham gia vụ này à?'' cô nhìn cậu Tổng với vẻ như vừa phát hiện ra kẻ phá bĩnh.

``Fumi nhờ anh thôi,'' cậu nhún vai. ``Chính em ấy là người quyết định ai lồng mà.''

``\dots Lần đầu tiên em thấy mình bị thay vai bởi một con mèo thật,'' Tamaki rên rỉ, gục đầu xuống bàn.

``Ừm,'' Suzuki vừa lật trang ghi chú vừa thêm vào, ``cơ mà Tamaki-nee vẫn có tố chất làm diễn viên trong hình mà.''

``\textbf{Chị là vật chủ ở đây đấy nhé!!}'' nàng Mèo hoang hét lên, giậm chân như một cô bé con bị lấy mất đồ chơi yêu thích. ``\textbf{Chị muốn mọi người nghe thấy tiếng kêu của chị, không còn phải con mèo dở hơi nào đó!!!}''

Fumi cười hiền hậu như thường lệ: ``Không sao. Lần sau vẫn có vai cho em mà. Vai\dots\ xem nào, khán giả chẳng hạn?''

``\textbf{Không!!!}''

Tiếng hét ``\textbf{Nyaaaaa!!!}'' vang vọng khắp cả tầng khiến cho nàng Cáo bật cười, còn anh em nhà Hikawa chỉ cười trừ khi bị cô bạn trêu chọc.

Kết quả cuối cùng là: ba nữ sinh trung học ấy, sau một trận cãi cọ không nhỏ, đã tự tìm cách giải quyết ổn thỏa với nhau. Riêng Kaede - kẻ ``trừng phạt'' trong vụ dorayaki bị cướp - thì bị cấp trên nhắc nhở nghiêm khắc vì hành vi phá hoại tài sản công ty.

Còn về phần ``nàng mèo hoang ngây thơ vô số tội'', Tamaki dù đã cố biện hộ rằng Kaede là người ``ra tay trước'' nên cô mới đáp trả, thì vẫn không nhận được sự cảm thông từ ban quản lý. Cộng thêm việc bị cô bạn thân của mình phản bội với việc để một con mèo thật lồng thay càng khiến cô uất ức hơn.

Và thế là, như một thói quen kỳ lạ đã lặp đi lặp lại, Tamaki lại một lần nữa thoát khỏi mọi trách nhiệm như chẳng có gì xảy ra, nhưng quyết định ``nghỉ chơi'' với cô bạn Cáo (không kéo dài được lâu sau khi được cho một lon cá). Còn ``trưởng Câu lạc bộ nữ sinh'' lại có thêm một lý do nữa để ghi tên Tamaki vào danh sách những ``mối thù truyền kiếp'' của đời mình.

``\dots \textit{Giá mà em ấy học được một phần tính nết trầm ổn từ Okayu-chan\dots}'' Kuon đứng khoanh tay nhìn xa xăm, khẽ lẩm bẩm. ``\textit{Mà, Tamaki-chan là fan của Okayu thật mà, không lẽ không học được gì từ thần tượng của mình sao?}''

\ct

Chiều muộn hôm đó, Kuon và Suzuki lặng lẽ quay trở về khu ký túc, khi bầu trời ngoài kia đã ngả màu chàm đậm. Con đường về lặng lẽ, không còn tiếng động nào ngoài những bước chân đều đặn và ánh đèn vàng êm dịu từ các cửa sổ văn phòng dần khép lại sau ngày dài.

``\vie{Tamaki-chan đúng là khó hiểu thật đấy,''} Kuon lầm bầm, tay cho vào túi áo khoác.

``\vie{Cô ấy sống bản năng như một con mèo hoang đúng nghĩa,''} cô em gái đáp, tay vẫn đan vào nhau sau lưng. ``\vie{Nhưng cũng bởi vì thế, chị ấy mới khiến người ta khó ghét được.''}

``\vie{Hừm. Anh thì không chắc là nên yêu quý một `con mèo' mà chuyên đi trộm bánh và phá văn phòng đâu. Nếu vậy có khi bố anh đã đuổi ra đường từ thuở nào rồi.''}

Cô mỉm cười. ``\vie{Vậy mà anh vẫn quay lại bình luận cho chương trình có `con mèo' đó đấy thôi.''}

``\vie{Vì Fumi nhờ. Anh cũng chẳng muốn đi làm gì.''}

``\vie{Vâng, vì Fumi nhờ thôi\~{},''} cô bé lặp lại, nhưng giọng có phần trêu chọc hơn.

Cả hai im lặng bước tiếp, ánh hoàng hôn phản chiếu lên nền gạch nhạt dần dưới chân họ. Gió chiều tà khẽ lướt qua, mang theo mùi hương cây lá dịu dàng sau cơn mưa sớm.

``\vie{Dù sao thì,''} Kuon dừng lại một chút trước cửa khu ký túc xá. ``\vie{Thật hiếm khi có một ngày mà mọi thứ kết thúc mà không có vụ nổ nào hay bất cứ thứ gì vượt quá mức cần thiết xảy ra kể từ khi anh đi dọn văn phòng hôm kia.''}

``\vie{Anh nói vậy như thể bình thường nó luôn phải có vậy.''}

``\vie{Thường là thế thật.''}

Suzuki cười khẽ. ``\vie{Vậy thì hôm nay nên được coi là đặc biệt.''}

Cánh cửa khẽ mở, cả hai bước vào trong sự yên tĩnh êm đềm, khép lại một ngày thứ sáu, nơi một nàng mèo hoang nghịch ngợm, ba cô gái trung học, và hai người bạn kể chuyện, đã cùng nhau tạo nên một chương trình hỗn loạn\dots\ nhưng đầy sắc màu.

\end{document}